% titel: Blatt 3, A7: Hess'scher Satz (Standard-Reaktionsenthalpie)
% Teil II: der durchgerechnete Weg.

\subsection*{Gegeben und gesucht}
Zielreaktion bei $T = \qty{298.15}{K}$:
\begin{equation*}
  \mathrm{N_2O_4(g)} + 2\,\mathrm{N_2H_4(l)} \rightarrow
  3\,\mathrm{N_2(g)} + 4\,\mathrm{H_2O(g)}.
\end{equation*}
Gegeben (Standard-bildungs-relevante Teilreaktionen):
\begin{align*}
  \mathrm{NH_3(g)} &\rightarrow \tfrac12\mathrm{N_2(g)} + \tfrac32\mathrm{H_2(g)}
    &\Delta H_1^\circ &= +46.19,\\
  \mathrm{N_2O_4(g)} &\rightarrow 2\,\mathrm{NO_2(g)}
    &\Delta H_2^\circ &= +58.04,\\
  2\,\mathrm{NH_3(g)} &\rightarrow \mathrm{N_2H_4(l)} + \mathrm{H_2(g)}
    &\Delta H_3^\circ &= +142.80,\\
  2\,\mathrm{NO_2(g)} + 4\,\mathrm{H_2(g)} &\rightarrow \mathrm{N_2(g)} + 4\,\mathrm{H_2O(g)}
    &\Delta H_4^\circ &= -1033.66.
\end{align*}
Gesucht: $\Delta_rH^\circ$ in kJ/mol.

\subsection*{Ansatz}
Zielreaktion aus den Teilreaktionen zusammensetzen durch Umkehren und
Vervielfachen, dann die Enthalpien addieren (Hess).

\subsection*{Rechnung}
\begin{enumerate}
  \item Gegeben (R4): $2\,\mathrm{NO_2} + 4\,\mathrm{H_2} \rightarrow
        \mathrm{N_2} + 4\,\mathrm{H_2O}$,
        $\Delta H_4^\circ = -1033.66$.
  \item R2 liefert $\mathrm{NO_2}$: $\mathrm{N_2O_4} \rightarrow
        2\,\mathrm{NO_2}$, $+58.04$.
  \item R3 umkehren und verdoppeln, um $2\,\mathrm{N_2H_4}$ als Edukt zu haben:
        $2\,\mathrm{N_2H_4} + 2\,\mathrm{H_2} \rightarrow
        4\,\mathrm{NH_3}$, $\Delta H = -2\cdot 142.80 = -285.60$.
  \item R1 vervierfachen, um $4\,\mathrm{NH_3}$ wieder abzubauen:
        $4\,\mathrm{NH_3} \rightarrow 2\,\mathrm{N_2} + 6\,\mathrm{H_2}$,
        $\Delta H = +4\cdot 46.19 = +184.76$.
\end{enumerate}
Addiert man (1)--(4), kürzen sich $2\,\mathrm{NO_2}$, $4\,\mathrm{NH_3}$ und
$6\,\mathrm{H_2}$ heraus:
\begin{align*}
  \Delta_rH^\circ &=
  \Delta H_4^\circ + \Delta H_2^\circ -2\,\Delta H_3^\circ +4\,\Delta H_1^\circ \\[2pt]
  &= (-1033.66) + 58.04 - 285.60 + 184.76 \\
  &= \qty{-1076.46}{kJ/mol}.
\end{align*}

\begin{ergebnis}
$\Delta_rH^\circ \approx \qty{-1076.46}{kJ/mol}$ für
$\mathrm{N_2O_4(g)} + 2\,\mathrm{N_2H_4(l)}
\rightarrow 3\,\mathrm{N_2(g)} + 4\,\mathrm{H_2O(g)}$.
\end{ergebnis}

\begin{kontrolle}
Probe der Teilbilanzen: R3 umgekehrt gibt $2\,\mathrm{N_2H_4}$, R4 liefert
$4\,\mathrm{H_2O}$ und (nach Kürzung) $1\,\mathrm{N_2}$,
R1' liefert $2\,\mathrm{N_2}$, zusammen $3\,\mathrm{N_2}$ --- passt zur
Zielreaktion. Die Reaktion ist stark exotherm (grosser negativer Wert),
typisch für ein Treibstoff-Paar ($\mathrm{N_2O_4}/\mathrm{N_2H_4}$).
\end{kontrolle}

\begin{merksatz}
Zielgleichung aus Teilgleichungen durch Umkehren/Vervielfachen zusammenbauen,
Zwischenprodukte kürzen und die (skalierten) $\Delta H$ aufsummieren.
\end{merksatz}
